% a01_H0rel part 7: assembly, normalisation-free form, the clause (R_0)
\subsection{Assembly, the normalisation-free form, and the clause \Rzero}\label{app:a01:assembly}

\begin{lemma}[field-independent constants cancel]\label{lem:a01:cnum}
Let $Z'(s)$, $I'(s)$ be $\Zk(s)$ and $\Ik j(s)$ computed with the effective action shifted by a field- and source-independent constant $c$ (a different $E_1$, a different normalisation of the measure, the domain-adapted $E_0$, $E_1$ of \cite[(3.35)]{K1}). If $e^{-a}I'\le Z'\le e^aI'+bZ'(0)$, then $e^{-a}\Ik j\le\Zk\le e^a\Ik j+b\Zk(0)$ with the same $a,b$.
\end{lemma}
\begin{proof}
$Z'=e^{-c}\Zk$ and $I'=e^{-c}\Ik j$ for every $s$; divide.
\end{proof}

\begin{proof}[Proof of \cref{prop:k01:H0rel}]
Fix $\eta_1\le\min(5^{-J_*},\eps_0^{\mathrm{Ba}})$ and smaller where the conditions on $\eta_1$ collected above require it (\cref{lem:a01:GC,lem:a01:final,prop:a01:legsd}, D7, D14), so that \cite[(3.54)]{Ba1} holds at every scale $\le j$ for $b_0\ge b_0^*$; the condition $j\ge J_*$ is Ba{\l}aban's $\eps\le\eps_0(\eta)=\eta5^{-J_*}$ and $j\le k-J_*$ is King's $\eta\le5^{-J_*}$, so that $j\le K_f\le k$. \Cref{prop:a01:lower,prop:a01:upper} give, in the normalisation of B6,
\[
e^{-a}\Ik j(s)\le\Zk(s)\le e^{a}\Ik j(s)+e^{-\frac12c_0'p(\eta)^2+E_+|\CJ|},\qquad\Zk(0)\ge e^{-E_-|\CJ|},
\]
hence the remainder is $\le e^{-\frac12c_0'p(\eta)^2+(E_++E_-)|\CJ|}\Zk(0)$. By \cref{lem:a01:cnum} the same holds for $\Zk$ and $\Ik j$ in the normalisation of \cref{ssec:k01:setting}. This is \eqref{eq:k01:H0rel} with $C\eta^\sigma$ in place of $C\eta^\sigma|\CJ|$ (the volume $|\CJ|=(2MJ)^3$ is absorbed in $C$, which may depend on $J$), $c=\frac12c_0'$ and $\rho(\eta)=e^{-cp(\eta)^2+C}$ with $C\ge(E_++E_-)|\CJ|$. The constants are independent of $k$ and $j$: of $j$ (i.e.\ of $\eta\le\eta_1$) because every step's constants are, and of $k$ (i.e.\ of $\vth$) because no step depends on the spacing except through the lattice regime, which \cref{lem:a01:DV,lem:a01:DVloc,lem:a01:A5} cover with the same constants; the dependence on $x$ enters through the source vertices only (B7, \cref{prop:a01:TD}). The clause \Rzero\ is \cref{thm:a01:R0} below.
\end{proof}

\begin{theorem}[the remainder is the large-field part of Ba{\l}aban's decomposition]\label{thm:a01:R0}
In the setting of \cref{prop:k01:H0rel}, let $P\defeq\prod_{i\le j}P^{(i)}$, where $P^{(i)}$ is the product, over all points, bonds and blocks of the scale-$i$ index sets of the all-small history (the whole cube's lattices at every scale), of the small-field indicators of the families of scale $i$ (Ba{\l}aban's six families, the face family, and at $i=j$ King's $\chi_j$ and $\lambda_j$). Then: (a) $P$ is a function on the joint variable space of Ba{\l}aban's block-spin representation with values in $\{0,1\}$, the same for every source $s$; (b) $\int P\,d\mu_{k,s}$ is the exact contribution of the all-small history at scale $j$, and $\int P\,d\mu_{k,0}\le e^{C\eta^\sigma}\Ik j(0)$; (c) $\int(1-P)\,d\mu_{k,0}\le e^{-\frac12c_0'p(\eta)^2+E_+|\CJ|}\le\rho(\eta)\Zk(0)$.
\end{theorem}
\begin{proof}
\emph{The joint measure.} Each block-spin step inserts $1=N_i^{-1}\int dB\,d\psi\exp(-\frac a2|B-QA|^2-\frac a2|\psi-Q(A)\phi|^2)$ \cite[(2.1)]{Ba2}, with $N_i$ field-independent; so $\Zk(s)=N^{-1}\int d\mu_{k,s}$ with $d\mu_{k,s}\defeq e^{-S}e^{X(s)}\prod_i[\text{kernels}]\,d(\text{all variables})$ a positive measure, and $N$ independent of $s$. \emph{The tested quantities are source-free:} they are fields, block fields and fluctuation fields (differences between a field and a background, \cite[(3.16)]{K1}), and backgrounds are linear maps of block fields built from the source-free quadratic forms; the source enters only the vertices \cite[(3.20)]{K1} and the factor $e^{X(s)}$. This gives (a). \emph{Decomposition at one scale.} Expanding $\prod[\chi(\text{small})+\chi(\text{large})]$ over the points, bonds and blocks of scale $i$ gives \cite[(2.4)--(2.6)]{Ba2}: $1=$ a sum over large-field tuples of products of indicators; the regrouping by $\Lambda_0$ and the partial resummation \cite[(2.9)--(2.10)]{Ba2} replace, for fixed $\Lambda_0\ne\Xl$, the indicator of ``the true tuple generates $\Lambda_0$'' by a sum over minimal admissible tuples of non-negative indicator products, which is $\ge$ the indicator it replaces; for $\Lambda_0=\Xl$ the only admissible tuple is the empty one and its term is exactly the product of all small-field indicators of that scale. \emph{Induction over scales.} In each history the next scale's decomposition is inserted where the RG continues; inside a term the translations are changes of variables with unit Jacobian, the Gaussian integrations are exact, and the removal of interactions from large-field regions and the large-field estimates \cite[(2.11)--(2.14)]{Ba2} are termwise upper bounds of a positive integrand; this is why \cite[(2.43)]{Ba2} is an inequality bounding $\Zk$ itself. \emph{Telescoping}, exact and pointwise, with $P_i\defeq\prod_{i'\le i}P^{(i')}$, $P_{-1}\defeq1$: $1-P=\sum_{i=0}^{j}P_{i-1}(1-P^{(i)})$, the $i$-th summand being the indicator of ``all tests of scales $<i$ pass everywhere, some test of scale $i$ fails''; by the identity-or-over-count at each scale, $P_{i-1}(1-P^{(i)})\le$ the sum of the weights $w_H$ of the histories $H$ whose first large-field block is at scale $i$, and summing over $i$, $1-P\le\sum_{H\ne H_0}w_H$, $H_0$ the all-small history (weight $P$). Integrate against the positive measure $d\mu_{k,0}$: $N^{-1}\int(1-P)d\mu_{k,0}\le\sum_{H\ne H_0}[\text{the (2.43)-term of }H\text{ continued to the final scale}]$, each history's exact contribution being bounded by its term because every operation after the insertion of the indicators is an identity or a termwise upper bound. The proof of \cref{prop:a01:upper} bounds exactly this sum by $e^{-\frac12c_0'p(\eta)^2+E_+|\CJ|}$, and $\Zk(0)\ge e^{-E_-|\CJ|}$ in the same normalisation (B6), which gives (c). (b): the all-small history's exact contribution is $\int P\,d\mu_{k,s}$, which the proof of \cref{prop:a01:upper} bounds by $e^{a}\Ik j(s)$. The normalisation-free form follows by \cref{lem:a01:cnum}. Positivity of $(1-P)\,d\mu_{k,0}$ holds because $P\le1$; whether \cite[(2.9)]{Ba2} is an identity or an over-count does not matter, only $1-P\le\sum_{H\ne H_0}w_H$ is used.
\end{proof}
